\begin{abstract}

To demystify the ``black box'' of audience voting in \textit{Dancing with the Stars}, we develop a comprehensive framework addressing the \textbf{constrained inverse problem} of reconstructing unobservable fan votes across 34 seasons.

\textbf{Vote Reconstruction (Problem 1): Two-Stage Bayesian Inference Engine.} We tackle the information bottleneck by constructing informative priors through a calibrated power-law model ($V_i \propto S_i^{1.2}$) fusing judge scores, celebrity attributes, and Google Trends. By employing SLSQP constrained optimization to align estimates with elimination outcomes, our model yields a \textbf{91.7\% prediction accuracy} overall (\textbf{99.5\% for the percentage method}). A median uncertainty measure of \textbf{CV = 0.117} validates the stability of our reconstruction.

\textbf{Voting Method Evaluation (Problem 2): Counterfactual Simulation.} Rigorous backtesting reveals that the rank-based method \textbf{amplifies fan influence by 12.8\%} relative to the percentage method in high-score density scenarios. While producer intervention mechanisms (judges selecting from the bottom two) reduce upset victories by 64\%, they simultaneously decrease social media engagement by 31\%, quantifying the trade-off between fairness and excitement.

\textbf{Mechanism Identification (Problem 3): GBDT-SHAP Analysis.} Leveraging Gradient Boosting Decision Trees with SHAP decomposition, we uncover divergent evaluation standards: partner experience drives \textbf{18.2\% of judge scores} (prioritizing technique), whereas celebrity industry background explains \textbf{23.6\% of fan vote variance} (favoring charisma).

\textbf{System Innovation (Problem 4): The ``Dramatic Arc'' Protocol.} We propose a phase-adaptive system that dynamically shifts weights (62\% judge-led to 68\% fan-led) to mimic narrative climax. Incorporating a controversy bonus mechanism, this protocol achieves a top-tier \textbf{composite score of 0.860}. It projects \textbf{47\% growth in engagement} while maintaining an optimal \textbf{11.1\% controversy rate}. Monte Carlo simulations confirm \textbf{89.3\% accuracy retention} under stochastic perturbations.

Our work delivers a \textbf{production-ready protocol} that reconciles competitive integrity with entertainment imperatives, providing quantified strategic value for reality television producers.

\begin{keywords}
Bayesian Inverse Modeling; Constrained Optimization; SHAP Explainability; Voting System Design; Fairness-Excitement Tradeoff
\end{keywords}

\end{abstract}

